% Fichier généré automatiquement par tools/generate_corrections.py.
% Source : CIN/CIN-01-Parametrage/09_RT_RSG/09_RT_RSG.tex
% Statut : complété (3/3 question(s)).
\begin{corrigebox}[Corrigé complété]
\CorrigeQuestion{1}
On associe au bâti le repère $\mathcal R_0=(O,\vec i_0,\vec j_0,\vec k_0)$.
            Le solide 1 roule sans glisser sur 0 : son orientation est repérée par
            $\theta(t)=(\vec i_0,\vec i_1)$ et son rayon est $R$. Le coulisseau 2 se
            déplace par rapport à 1 suivant $\vec i_1$ avec
            $\overrightarrow{AB}=\lambda(t)\vec i_1$. La condition de roulement au point
            $I$ impose $\vec V(I,1/0)=\vec0$ et relie l'abscisse du centre $A$ à
            $\theta$ par $x_A=x_A(0)-R\theta$ (signe adapté à la convention). Ainsi
            $\vec\Omega_{1/0}=\dot\theta\vec k_0$ et
            $\vec\Omega_{2/1}=\vec0$.
\CorrigeQuestion{2}
\CorrigeSource{Réponse reprise d'un exercice homonyme de la banque ; la provenance exacte est indiquée dans le commentaire du fichier.}
% Réponse reprise de : CIN/CIN-02-VitesseAcceleration/09_RT_RSG/09_RT_RSG.tex
$\torseurcin{V}{2}{0}=\torseurl{\thetap \vk{0}}{ \lambdap\vi{1} + \thetap\left(\lambda(t)\vj{1}-R\vi{0} \right)}{B}$.
\CorrigeQuestion{3}
\CorrigeSource{Réponse reprise d'un exercice homonyme de la banque ; la provenance exacte est indiquée dans le commentaire du fichier.}
% Réponse reprise de : CIN/CIN-02-VitesseAcceleration/09_RT_RSG/09_RT_RSG.tex
$\vectg{B}{2}{0} = \deriv{\vectv{B}{2}{0}}{\rep{0}}$
$ = \lambdapp(t)\vi{1} +\lambdap(t)\thetap\vj{1} 
+ \thetapp(t)\left(\lambda(t)\vj{1}-R\vi{0} \right)
+ \thetap(t)\left(\lambdap(t)\vj{1}-\lambda(t)\thetap\vi{1} \right)
$.

\begin{enumerate}
\item $\vectv{B}{2}{0} = \lambdap\vi{1} + \thetap\left(\lambda(t)\vj{1}-R\vi{0} \right)$.
\item $\torseurcin{V}{2}{0}=\torseurl{\thetap \vk{0}}{ \lambdap\vi{1} + \thetap\left(\lambda(t)\vj{1}-R\vi{0} \right)}{B}$.
\item $\vectg{B}{2}{0}  = \lambdapp(t)\vi{1} +\lambdap(t)\thetap\vj{1} 
+ \thetapp(t)\left(\lambda(t)\vj{1}-R\vi{0} \right)
+ \thetap(t)\left(\lambdap(t)\vj{1}-\lambda(t)\thetap\vi{1} \right)
$.
\end{enumerate}
\end{corrigebox}
